\documentclass{article}

\usepackage{amsmath,amssymb}
\usepackage{xurl}
\usepackage[colorlinks=true,linkcolor=blue,citecolor=blue,urlcolor=blue]{hyperref}

% Shared commands for notes in docs/paper-gaps/.

\DeclareUnicodeCharacter{2080}{\ifmmode{}_{0}\else\textsubscript{0}\fi}
\DeclareUnicodeCharacter{2081}{\ifmmode{}_{1}\else\textsubscript{1}\fi}
\DeclareUnicodeCharacter{2082}{\ifmmode{}_{2}\else\textsubscript{2}\fi}
\DeclareUnicodeCharacter{2099}{\ifmmode{}_{n}\else\textsubscript{n}\fi}

\newcommand{\Lean}{\textsc{Lean}}
\ExplSyntaxOn
\NewDocumentCommand{\leanid}{m}{
  \ifmmode
    \textsf{\detokenize{#1}}
  \else
    \group_begin:
    \urlstyle{sf}
    \tl_set:Nn \l_tmpa_tl {#1}
    \tl_replace_all:Nnn \l_tmpa_tl {\_} {_}
    \exp_args:NV \path \l_tmpa_tl
    \group_end:
  \fi
}
\ExplSyntaxOff
\providecommand{\tr}{\operatorname{tr}}
\newcommand{\tnleanIssueBase}{https://github.com/LionSR/TNLean/issues/}
\newcommand{\tnleanPrBase}{https://github.com/LionSR/TNLean/pull/}
\newcommand{\tnleanDocBase}{https://sirui-lu.com/TNLean/docs/}
\newcommand{\ghissue}[1]{\href{\tnleanIssueBase#1}{GitHub issue \##1}}
\newcommand{\ghpr}[1]{\href{\tnleanPrBase#1}{PR \##1}}
\newcommand{\doclink}[1]{\href{\tnleanDocBase#1.html}{\path{#1}}}

% Dirac notation and the identity operator, as in blueprint/src/macros/common.tex.
% \providecommand keeps note-local definitions authoritative.
\usepackage{dsfont}
\providecommand{\ket}[1]{|#1\rangle}
\providecommand{\bra}[1]{\langle #1|}
\providecommand{\braket}[2]{\langle #1 | #2 \rangle}
\providecommand{\ketbra}[2]{|#1\rangle\!\langle #2|}
\providecommand{\Id}{\mathds{1}}
\providecommand{\one}{\mathds{1}}
\providecommand{\C}{\mathbb{C}}
\providecommand{\MN}[1]{M_{#1}(\C)}
\providecommand{\MD}{M_D(\C)}

% Verdict marker: \gapnote{<kind>}{<status>}. Kind is the policy
% classification (clarification, local-correction, scope-restriction,
% unfaithful, false-source, open-gap); status is open, wip (elimination
% underway), resolved, or historical. Machine-read by the site index;
% prints a verdict line.
\newcommand{\gapnote}[2]{\par\noindent\textbf{Verdict:} #1 (#2).\par}


\title{A Direct Cross-Section Replacing the Slice Theorem in the Continuous
    Standard Form}
\author{}
\date{2026-11-16}

\begin{document}
\maketitle
\gapnote{clarification}{wip}

\section{The source step}

This note concerns the final step in the proof of
Corollary~\texttt{coruvSF} of Ref.~\cite[lines~942--947]{Cirac2017MPU},
the continuous choice of standard form for a path of matrix product unitaries.
At that point the proof has shown that the blocked gate $u_{12}(x)$, read
modulo a local-unitary gauge, depends continuously on the path parameter $x$:
for nearby $x$ and $y$ there are matrices $r_1(x,y)$ and $s_2(x,y)$ acting on
the two tensor factors with
\begin{align}
    \bigl\|u_{12}(x)-\bigl(r_1(x,y)\otimes s_2(x,y)\bigr)\,u_{12}(y)\bigr\|_1
        &\leq C\varepsilon.
    \label{eq:mpu_standard_form_path_cross_section_local_gauge}
\end{align}
To upgrade this to a globally continuous gate, the source invokes a classical
result of Montgomery and Yang on the existence of a slice for the action of the
compact group of local unitaries $U(\ell)\otimes U(r)$. The slice theorem
produces a continuous local-unitary gauge $\hat r(x)\otimes\hat s(x)$ such that
\begin{align}
    \hat u(x)
        &=\bigl(\hat r(x)\otimes\hat s(x)\bigr)\,u(x)
    \label{eq:mpu_standard_form_path_cross_section_slice_gauge}
\end{align}
is continuous in $x$.

\section{The formal route}

The slice theorem of Montgomery and Yang is a statement about the orbit space
of a smooth compact-group action and carries substantial differential-topology
infrastructure that the formal development does not have. The formalization
replaces it with a direct construction of a continuous cross-section for the
local-unitary gauge, built from elementary linear algebra on the realignment of
a bipartite matrix.

The realignment pairs the two indices of each tensor factor. For a matrix $M$
from $\mathbb C^{m'}\otimes\mathbb C^{n'}$ to
$\mathbb C^{m}\otimes\mathbb C^{n}$, its realignment $R(M)$ is the matrix with
rows indexed by pairs $(i,i')$, columns by pairs $(k,k')$, and entries
\begin{align}
    R(M)_{(i,i'),(k,k')}
        &=M_{(i,k),(i',k')}.
    \label{eq:mpu_standard_form_path_cross_section_realign_entry}
\end{align}
Two identities make $R$ the right tool for isolating the Kronecker factors of a
gate. First, on a Kronecker product it returns the outer product of the entry
vectors,
\begin{align}
    R(A\otimes B)
        &=(\operatorname{vec}A)\,(\operatorname{vec}B)^{\mathsf T}.
    \label{eq:mpu_standard_form_path_cross_section_realign_kronecker}
\end{align}
Second, it is equivariant under left multiplication by a local operator
$h\otimes k$,
\begin{align}
    R\bigl((h\otimes k)\,M\bigr)
        &=(h\otimes\mathbb 1)\,R(M)\,(k\otimes\mathbb 1)^{\mathsf T}.
    \label{eq:mpu_standard_form_path_cross_section_realign_equivariance}
\end{align}
Together with the rank characterization
\begin{align}
    \operatorname{rank}R(M)=1
        &\iff M=A\otimes B,\quad A\neq 0,\ B\neq 0,
    \label{eq:mpu_standard_form_path_cross_section_rank_one}
\end{align}
these identities convert the approximate local-gauge statement
(\ref{eq:mpu_standard_form_path_cross_section_local_gauge}) into a statement
about the leading singular vectors of $R\bigl(u_{12}(x)u_{12}(y)^\dagger\bigr)$,
which vary continuously with $(x,y)$ and supply the continuous factors
$\hat r$ and $\hat s$ of
(\ref{eq:mpu_standard_form_path_cross_section_slice_gauge}) directly, without
appeal to a slice theorem.

The realignment and its three identities are the formalized first step of this
construction.\footnote{Formal declarations:
    \leanid{Matrix.realign},
    \leanid{Matrix.realign_kronecker},
    \leanid{Matrix.realign_kronecker_mul}, and
    \leanid{Matrix.rank_realign_eq_one_iff} in
    \path{TNLean/Algebra/MatrixRealignment.lean}.}
Equation~(\ref{eq:mpu_standard_form_path_cross_section_realign_entry}) is
\leanid{Matrix.realign};
(\ref{eq:mpu_standard_form_path_cross_section_realign_kronecker}) is
\leanid{Matrix.realign_kronecker};
(\ref{eq:mpu_standard_form_path_cross_section_realign_equivariance}) is
\leanid{Matrix.realign_kronecker_mul}; and
(\ref{eq:mpu_standard_form_path_cross_section_rank_one}) is
\leanid{Matrix.rank_realign_eq_one_iff}.

\section{Conclusion}

The replacement is a change of proof route, not of statement. The conclusion of
Corollary~\texttt{coruvSF} is unchanged: a continuous path of standard forms
exists for every continuous path of tensors. Where the source cites the
Montgomery--Yang slice theorem for the compact local-unitary action, the
formalization constructs the continuous gauge explicitly from the realignment
and its singular-value data. The remaining steps of the construction, and the
corollary itself, are not yet formalized; the corollary is marked
\texttt{\textbackslash notready} in the blueprint and this note is a
work-in-progress record of the planned route.

\bibliographystyle{alpha}
\bibliography{references}

\end{document}
